\section{From the valley profiles to one polynomial}\label{sec:assembly}
\subsection{Separation of the radial scales}\label{subsec:inductive-valleys}
De Maesschalck and Huzak preserve earlier simple zeros while adding
powers at a larger radius \cite[\S3.1]{DMH}. Here each new radius must
also preserve the ordering of the critical levels and three finite lists
of signs. Fix $D=258$, $L\ge1$ and
$j\in\{0,1,2,3\}$. For an odd initial integer $n_0$ put
$n_i=n_0+3i$ for $0\le i\le L$, $m_0=129n_0$, $d_\nu=Dn_\nu$, $\ell_0=2j$, and $\ell_\nu=0$
for $\nu>0$. The support after the last fold is
\begin{equation}\label{eq:support-general}
\begin{aligned}
 S_{L,j}(n_0)={}&\{m_0,d_0-1,d_0\}\cup\{d_0-1+2i:1\le i\le j\}\\
 &{}\cup\bigcup_{\nu=0}^{L-1}
 \{d_\nu+\ell_\nu+D-1,d_\nu+2D,d_\nu+3D-1,d_\nu+3D\}\\
 &{}\cup\{d_L+1,d_L+3,d_L+8\}.
\end{aligned}
\end{equation}
Take $N_*$ to be the maximum of $5$, the seed threshold and the four
local-profile thresholds in Lemmas~\ref{lem:compact-seed}
and~\ref{lem:valley-local}, expressed in the variable $n_0$.
The block using $d_i$ has parameter $n_i$, for $0\le i<L$. The parameters increase with $i$.
Thus every odd $n_0\ge N_*$ works for every finite $L$ and every $j$.
For fixed $(n_0,L,j)$ we choose the block radii successively for the
finite tests, then dilate the assembled profile and balance its tails.
Only after the physical shift and test heights are fixed do we choose
the outer radius, solve the three matches and regularize the returns.
For a fixed such $n_0$, the largest exponent in \eqref{eq:support-general}
is $d_L+8=258n_0+774L+8$. Since $k=4L+6+j$, these supports have degree
$\tfrac{387}{2}k+O_{n_0}(1)$ as $k$ increases, uniformly over the four
values of $j$.

\begin{lemma}[Multiscale assembly of $j$ prefix terms and $L$ positive valleys]\label{lem:inductive-valleys}
Fix $D=258$, $L\ge 1$, and $j\in\{0,1,2,3\}$, and choose an odd integer $n_0$ sufficiently large that
Lemma~\ref{lem:compact-seed} holds at $m_0=129n_0$ and Lemma~\ref{lem:valley-local} holds for
$n_i=n_0+3i$ ($0\le i<L$). (In particular, for $L=1$ and $j=0$, $n_0=257$ suffices by
Section~\ref{subsec:ten-power}.) Then there exist polynomials $E(r),Q(r)$ with collected supports
\begin{align*}
  \supp(E) &= \{d_0\} \cup \bigcup_{\nu=0}^{L-1}\{d_0+3D\nu+2D,\, d_0+3D\nu+3D\}, \\
  \supp(Q) &= \{d_0-1\} \cup \{d_0-1+2i:1\le i\le j\} \cup \{d_0+2j+D-1,\, d_0+3D-1\} \\
           &\quad \cup \bigcup_{\nu=1}^{L-1}\{d_0+3D\nu+D-1,\, d_0+3D\nu+3D-1\},
\end{align*}
such that $E$ has positive leading coefficient at $d_{\mathrm{top}}=d_L=d_0+3DL$, $Q$ has leading
coefficient of sign $\omega=(-1)^j$ at $d_{\mathrm{top}}-1$, $E>0$ on $(0,\infty)$ satisfies the
ordered positive-valley hypotheses \eqref{eq:valley-ordering} with $L$ crests and $L$ minima, and the
variations $V_o(h)$, $V_i(h)$, and $J(H)$ of Lemma~\ref{lem:valley-rates} have respectively $6L+j$,
$6L+j$, and $4L+j$ strict sign changes on ordered finite test lists. Before reflection, each list
starts strictly positive at small $h$ or $H$ and ends with sign $\omega=(-1)^j$ at large $h$ or $H$.
\end{lemma}

\begin{proof}
Begin with $E_0(r)=r^{d_0}$ and $Q_0(r)=r^{d_0-1}+\sum_{i=1}^j b^{(i)} r^{d_0-1+2i}$. While
$E_0(r)=r^{d_0}$ ($f_0(r)=d_0 r^{e_0}$, $e_0=d_0-2$), an odd term $b_p r^p$ with $d_0-1\le p\le
d_0+5<2d_0-2$ contributes
\[
  b_p \frac{p-e_0}{p+e_0} r^p \quad \text{to } \frac{U(r)}{f_0(r)}, \qquad b_p H^{(p+2)/d_0 - 2}
  \int_0^\infty \frac{t^{p+1}}{(1+t^{d_0})^2}\dd t \quad \text{to } J(H),
\]
both with strictly positive multipliers of $b_p$. Choosing $b^{(1)},b^{(2)},\dots,b^{(j)}$
successively small in magnitude with $\sgn b^{(i)}=(-1)^i$ preserves the earlier signs at both
positive and negative heights and adds a larger test in $r$ and $H$ dominated by $b^{(i)}$, yielding $j$ sign changes
(starting $+$ and ending $(-1)^j$) on each of $V_o,V_i,J$, with leading odd sign $\beta_0=b^{(j)}$
($b^{(0)}=1$) equal to $(-1)^j$.

Now induct over $\nu=1,\dots,L$. Suppose after step $\nu-1$ that $(E,Q)$ has leading terms $\alpha
r^d$ ($\alpha>0$, $d=d_{\nu-1}=Dn_{\nu-1}$) and $\beta r^{d-1+l}$ ($\sgn\beta=(-1)^j$), where $l=2j$ for
$\nu=1$ and $l=0$ for $\nu\ge 2$. Set $\gamma=l/D$, $\eta=(2/3)^{n_{\nu-1}}$, $q=(3/4)^{\lfloor
n_{\nu-1}/2\rfloor}$, $z_S=(r/S)^D$, and define at a large radial scale $S>1$:
\[
  E_S(r) = E(r) + \alpha r^d\bigl[C_\eta(z_S) - 1\bigr], \qquad Q_S(r) = Q(r) + \beta r^{d-1+l}\bigl[B_\gamma(z_S) - qz_S - 1\bigr].
\]
Because $C_\eta(z)-1=-(3-\eta)z^2+2z^3$ and $B_\gamma(z)-qz-1=-(9/8+q)z+\frac{1}{8}z^{3-\gamma}$ have
zero constant terms and nonzero higher coefficients, all old coefficients remain untouched, the four
added exponents $d+2D$, $d+3D$, $d+l+D-1$, and $d+3D-1$ exceed all previous exponents in their parity
with nonzero physical coefficients $-(3-\eta)\alpha/S^{2D}$, $2\alpha/S^{3D}$, $-(9/8+q)\beta/S^D$,
and $\beta/(8S^{3D-l})$, the new physical top even coefficient is $2\alpha/S^{3D}>0$ at $d+3D$, and
the new physical top odd coefficient is $\beta/(8S^{3D-l})$ (preserving sign $(-1)^j$) at $d+3D-1$
(with normalized model leading coefficients $2\alpha$ and $\frac{1}{8}\beta$ in $t=r/S$).

As $S\to+\infty$, $(E_S,Q_S)\to(E,Q)$ in $C^2$ on compact intervals of $[0,\infty)$, while at $r=St$,
$(\alpha^{-1}S^{-d}E_S(St),\beta^{-1}S^{-(d-1+l)}Q_S(St))$ converges in $C^2$ on compact subsets of
$(0,\infty)$ to the isolated model pair $(E_{\mathrm{loc}}(t),Q_{\mathrm{loc}}(t))$ of
Lemma~\ref{lem:valley-local}. On $r\ge 1$, subleading powers of $(E,Q)$ contribute $O(r^{-\kappa})$
($\kappa\ge 2$):
\[
  \frac{E_S(r)}{\alpha r^d} = C_\eta(z_S) + O(r^{-\kappa}), \qquad \frac{E_S'(r)}{\alpha d r^{d-1}} =
  C_\eta(z_S) + \frac{z_S}{n_{\nu-1}}C_\eta'(z_S) + O(r^{-\kappa}).
\]
Since $C_\eta(z)\ge c_*(n_{\nu-1})(1+z^3)>0$ on $z\ge 0$ and $C_\eta(z)+zC_\eta'(z)/n_{\nu-1}$ has nonzero
constant and cubic coefficients and only the two simple positive roots $z_p,z_c$, no critical points
can appear in the intermediate or far regimes $r\to\infty$ with $z_S\to 0$ or $z_S\to\infty$. Near $0$
divide by $r^{d_0}$. Compact $C^1$ convergence therefore implies $E_S>0$ on $(0,\infty)$, all $\nu-1$
old simple crest--valley pairs persist, and exactly one new crest $p_{\nu,S}\sim S z_p^{1/D}$ and
valley $c_{\nu,S}\sim S z_c^{1/D}$ appear, with minimum value $E_S(c_{\nu,S})\sim \alpha S^d
m_{\mathrm{loc}}\to+\infty$ exceeding all old critical values and old test heights.

Consequently, all old positive-height tests lie below $E_S(c_{\nu,S})$ and their variations $V_o,V_i$
converge to their old nonzero values. At new positive test heights $\alpha S^d h_{\mathrm{loc}}$,
dividing $U_S(St)$, $f_S(a_{\nu,S})Q_S(c_{\nu,S})$, and $U_S(c_{\nu,S})$ by $\alpha|\beta| S^{2d-3+l}$
sends the bounded old sum $C_{\nu-1}^*$ to $0$ and converges to $\sgn(\beta)$ times the isolated model
variations of Lemma~\ref{lem:valley-local}, whose first sign $\sgn\beta=(-1)^j$ matches the final sign
of the old list and adds $6$ new sign changes to both $V_o$ and $V_i$.

For $J_S(H)$, on $r\ge K_1$ and $S\ge 1$ one has $E_S(r)\ge c_1 r^d(1+z_S^3)$ and $|Q_S(r)|\le C_1
r^{d-1+l}(1+z_S^3)$, yielding the $S$-uniform integrable tail majorant $C_2 r^{l-d}\le C_2 r^{6-d_0}$
at each fixed old depth $H>0$. At each of the five new depths $H=\alpha S^d H_{\mathrm{loc}}$,
rescaling $r=St$ bounds the normalized denominator below by $H_{\mathrm{loc}}>0$ on $[0,1]$ and by
$c_1 t^d(1+t^{3D})$ on $[1,\infty)$, so dominated convergence yields $\sgn\beta=(-1)^j$ times the five
alternating signs $+,-,+,-,+$ of Lemma~\ref{lem:valley-local}, joining the old final sign $(-1)^j$
without loss and adding $4$ new sign changes. Choosing $S$ large enough for all three finite lists
simultaneously completes the induction.
\end{proof}
